%% CV - Giovanny Garcia
%% Tailored for Software Engineering Intern, UpAhead
%% Compile with: lualatex main_upahead.tex

\documentclass[11pt,a4paper,sans]{moderncv}
\moderncvstyle{banking}
\moderncvcolor{blue}

% Force both first and last name AND section headings to render in moderncv
% blue (color1). Default banking on lualatex+MiKTeX leaves these black, which
% looks inconsistent with the rest of the blue accent scheme.
\renewcommand*{\firstnamestyle}[1]{{\fontsize{34}{36}\bfseries\upshape\color{color1}#1}}
\renewcommand*{\lastnamestyle}[1]{{\fontsize{34}{36}\bfseries\upshape\color{color1}#1}}
\renewcommand*{\sectionstyle}[1]{{\sectionfont\color{color1}#1}}

\usepackage[utf8]{inputenc}
\usepackage{needspace}
\usepackage{hyperref}
\hypersetup{
    colorlinks=true,
    linkcolor=blue,
    filecolor=magenta,
    urlcolor=blue,
    pdftitle={Giovanny Garcia - CV},
    pdfpagemode=FullScreen,
}
\usepackage[scale=0.80]{geometry}
\usepackage{import}

% personal data
\name{Giovanny}{Garcia}
\address{Austin, Texas}{}{}
\email{giovanny.garcia2@g.austincc.edu}
\extrainfo{\href{https://www.linkedin.com/in/giovanny-garcia-07ab24322}{LinkedIn}, \href{https://github.com/giovanny-garcia}{GitHub}}

\begin{document}

\makecvtitle

% ============================================================
%     PROFILE STATEMENT
% ============================================================

\vspace{6pt}
\small{Computer science student at Austin Community College and daily Blackboard Ultra user. Ships TypeScript and Node.js features with slash-command UX, SQLite persistence, and a cache-first API design so the product stays useful when a third-party quota is tight. Independent React practice plus C++ and C\# from the CS toolkit. Uses Claude Code to multiply output on scoping, debugging, and iteration while keeping the architecture decisions. Next 12 weeks: React on the web app, React Native ramped on the live iOS app, aimed at usage on a semester view students already need.}

% ============================================================
%     CORE COMPETENCIES
% ============================================================

\section{Core Competencies}
\vspace{1pt}
\begin{itemize}

\item \textbf{TypeScript \& Node.js}: Ships typed backend features with discord.js, REST clients, structured error handling, and local persistence.

\item \textbf{Front-end (React)}: Component-based UI practice. Ready to contribute on the web app (React) and ramp on React Native for the live iOS app.

\item \textbf{SQL \& data modeling}: SQLite schemas for settings, caches, deduplication, and quota tracking so user-facing commands stay fast.

\item \textbf{Product-minded shipping}: Designs around real constraints (API quotas, reminder windows, per-server settings) so the product stays used, not just so ``it works on my machine.''

\item \textbf{AI-assisted development}: Uses Claude Code to multiply output on implementation, debugging, and refactors while keeping the architecture decisions.

\item \textbf{Git \& ownership}: Public GitHub history on a project I scoped, built, and still own end to end.

\end{itemize}

% ============================================================
%     INDEPENDENT PROJECTS
% ============================================================

\section{Independent Projects}
\vspace{3pt}
\begin{itemize}

\needspace{5\baselineskip}
\item{\cventry{2026}{CS2 Discord Bot}{\href{https://github.com/giovanny-garcia/hltv-discord-bot}{github.com/giovanny-garcia/hltv-discord-bot}}{TypeScript, Node.js, SQLite}{}{\vspace{1pt}
Personal product that announces upcoming Counter-Strike 2 matches and lets a server browse schedules without burning a scarce third-party API quota.
\begin{itemize}
    \item Built a cache-first architecture: slash commands such as \texttt{/matches}, \texttt{/results}, and \texttt{/events} read SQLite, while \texttt{/sync} is the only path that calls the GGScore REST API, so browsing still works when the daily quota is gone.
    \item Enforced a 3-request/day free-tier budget with quota tracking, typed API errors, and a \texttt{/quota} command so operators can see remaining calls.
    \item Added per-guild subscriptions, reminder windows, and a poll loop that posts embeds for new or starting-soon matches without extra API traffic.
    \item Modeled Discord as the client: slash-command UX, permission-gated admin commands, and settings that persist across restarts.
    \item Owned the full loop: TypeScript toolchain (tsx, tsc), environment config, SQLite via better-sqlite3, and a public README that documents the free-tier strategy.
\end{itemize}}}

\end{itemize}

% ============================================================
%     EDUCATION
% ============================================================

\section{Education}
\vspace{1pt}
\begin{itemize}

\needspace{5\baselineskip}
\item{\cventry{2022--2025}{Associate of Science in Computer Science}{Austin Community College}{Austin, Texas}{}{\vspace{1pt}
Transfer-oriented CS coursework covering programming fundamentals, object-oriented design, and data structures (C++ and C\#). JavaScript, TypeScript, and React are independent practice on top of that base. Daily student user of ACC's Blackboard Ultra LMS: due dates, grades, slides, and attendance policies split across courses, which is the exact problem UpAhead's LMS sync and syllabus extraction consolidate.
}}

\end{itemize}

% ============================================================
%     TECHNICAL SKILLS
% ============================================================

\section{Technical Skills}
\vspace{1pt}
\begin{itemize}

\item \textbf{Languages}: TypeScript, JavaScript, C++, C\#, SQL.

\item \textbf{Web \& app}: React, Node.js, REST APIs. React Native is a gap I would close on the job against the live iOS app, using React and TypeScript as the ramp.

\item \textbf{Data \& tools}: SQLite, Git, npm, Discord/slash-command UX, Claude Code.

\item \textbf{Practices}: quota-aware API clients, caching, structured errors, README-driven setup, iterating from real user constraints.

\end{itemize}

% ============================================================
%     ADDITIONAL BACKGROUND
% ============================================================

\section{Additional Background}
\vspace{1pt}
\begin{itemize}

\item \textbf{Game development}: Independent C\# practice listed alongside the CS coursework.

\item \textbf{How I work}: I pick a user problem, constrain the design (quotas, permissions, reminders), then ship a vertical slice and iterate. I already own scope end to end on my own project. Ready to do that with the founder on a small team, without waiting for a ticket queue.

\item \textbf{Work setup}: Austin, Texas. US-based. Available about 20 hours a week for the remote internship window (24 August 2026 -- 16 November 2026).

\end{itemize}

% ============================================================
%     REFERENCES
% ============================================================

\section{References}
\vspace{1pt}
\begin{itemize}
\item Available upon request.
\end{itemize}

\end{document}
